\documentclass[a4paper]{article}
\usepackage[margin=.8in]{geometry}
\usepackage{amsmath}
\usepackage{enumitem}
\usepackage[cache=true]{minted}
\usepackage{amsfonts}
\usepackage{amssymb}
\usepackage{mathtools}
\input{../../macros}

\usepackage{fancyhdr}
\pagestyle{fancy}
\fancyhf{}
\lhead{Name: }
\chead{Graded quiz 1 -- Floating-point arithmetic}
\rhead{\today}

\begin{document}

\pagestyle{empty}
\thispagestyle{fancy}

True or false? (unless otherwise specified.)

\begin{center}
    \textbf{Answer grid}\par\smallskip
    Circle one answer for each question.

    \smallskip
    \normalsize
    \renewcommand{\arraystretch}{1.35}
    \begin{tabular}{|c|*{17}{c|}}
        \hline
        & 1 & 2 & 3 & 4 & 5 & 6 & 7 & 8 & 9 & 10 & 11 & 12 & 13 & 14 & 15 & B1 & B2 \\
        \hline
        True  & $\square$ & $\square$ & $\square$ & $\square$ & $\square$ & $\square$ & $\square$ & $\square$ & $\square$ & $\square$ & $\square$ & $\square$ & $\square$ & $\square$ & $\square$ & $\square$ & $\square$ \\
        False & $\square$ & $\square$ & $\square$ & $\square$ & $\square$ & $\square$ & $\square$ & $\square$ & $\square$ & $\square$ & $\square$ & $\square$ & $\square$ & $\square$ & $\square$ & $\square$ & $\square$ \\
        \hline
    \end{tabular}
\end{center}

\begin{enumerate}

    \item
        The binary number $(101.101)_2$ is equal to the decimal number $5.625$.

    \item
        A real number has a finite representation in base 2 if and only if it is a
        rational number.

    \item
        The binary representation of $\frac{1}{3}$ is $(0.\overline{01})_2$.

    \item
        Let $(\placeholder)_2$ denote binary representation.
        Then
        \[
            (1000000000)_2 \times (0.001)_2 = (1000000)_2.
        \]

    \item
        It holds that
        \[
            (0.111111)_2 + (0.000001)_2 = (1.001)_2.
        \]

    \item
        It holds that $(0.\overline{1100})_2 = \frac{4}{5}$.

    \item
        Let $\varepsilon_{32}$ denote the machine epsilon for the \julia{Float32} format,
        i.e.\ \julia{eps(Float32)}.
        Then the number~3 is representable exactly in this format,
        and the next representable number in the \julia{Float32} format is
        $3 + 2\varepsilon_{32}$.

    \item
        In Julia, \julia{Float64(3.25) == Float32(3.25)} evaluates to \julia{true}.

    \item
        The spacing (in absolute value) between successive half-precision (\julia{Float16})
        floating point numbers is always greater than or equal to \julia{eps(Float16)}.

    \item
        In Julia, \julia{Float64(.6) == Float32(.6)} evaluates to \julia{true}.

    \item
        There exists $n \in \nat$ such that the machine epsilon for the \julia{Float64}
        format equals $2^{-n}$.
        In other words, the machine epsilon is an exact negative power of two.

    \item
        In Julia, it holds that \julia{eps()/3 + eps()/3 + 1 > 1} evaluates to \julia{true}.

    \item
        For any real number $x$ such that $x \in \floating_{64}$,
        it holds that $2x \in \floating_{64}$.

    \item
        Infinitely many distinct real numbers can be represented exactly in the~\julia{Float64}
        format, but only finitely many can be represented exactly in the \julia{Float32} format.

    \item
        Let $x$ and $y$ be two numbers in $\floating_{64}$.
        The result of the machine multiplication $x \mtimes y$ is sometimes exact and sometimes not,
        depending on the values of $x$ and $y$.

    \item
        \textbf{(Bonus)}
        Does the Julia expression \julia{sqrt(1 + 2eps()) == 1} evaluate to \julia{true}
        or \julia{false}?

        \textit{Explain briefly:}
        \vspace{1.5cm}

    \item
        \textbf{(Bonus)}
        Does the Julia expression \julia{log2(2 + eps()) == 1} evaluate to \julia{true}
        or~\julia{false}?

        \textit{Explain briefly:}
        \vspace{1.5cm}

\end{enumerate}

\ifdefined\showsolutions
\newpage
\lhead{Instructor copy}
\section*{Answer key}
\noindent Questions 1--15: 1 point for a correct answer, 0 otherwise.

\begin{enumerate}[itemsep=8pt]
    \item \textbf{True.}
        $(101.101)_2 = 4 + 1 + 1/2 + 1/8 = 5.625$.

    \item \textbf{False.}
        A binary expansion is finite if and only if the number is a
        dyadic rational, i.e.\ of the form $m/2^k$ for integers
        $m$ and $k$. Most rational numbers (e.g.\ $1/3$) require an
        infinite expansion.

    \item \textbf{True.}
        $1/3 = 0.010101\dotsc_2$, i.e.\ $(0.\overline{01})_2$.

    \item \textbf{True.}
        $(0.001)_2 = 2^{-3}$, so
        $2^9 \cdot 2^{-3} = 2^6 = (1000000)_2$.

    \item \textbf{False.}
        The last carry propagates through all the bits:
        $(0.111111)_2 + (0.000001)_2 = (1.000000)_2 = 1$, not
        $(1.001)_2$.

    \item \textbf{True.}
        The repeating block is $1100$, whose value is $3/4$, so
        $S = (3/4)/(1 - 1/16) = (3/4) \cdot (16/15) = 4/5$.

    \item \textbf{True.}
        In \julia{Float32} the spacing near 3 is
        $2 \cdot \varepsilon_{32}$ (the exponent is 1), so the next
        representable number is $3 + 2\varepsilon_{32}$.

    \item \textbf{True.}
        $3.25 = 13/4$ is a dyadic rational, so it is represented exactly
        in both formats and converting one to the other is exact.

    \item \textbf{False.}
        For numbers below 1, e.g.\ between $1/2$ and 1, the spacing in
        \julia{Float16} is $2^{-11}$, which is smaller than
        \julia{eps(Float16) = 2^{-10}}.

    \item \textbf{False.}
        The decimal number $0.6$ is not exactly representable in binary,
        and its nearest \julia{Float64} and \julia{Float32} values differ.

    \item \textbf{True.}
        \julia{eps(Float64)} is $2^{-52} = \left(\frac{1}{2}\right)^{52}$,
        an exact (negative) power of two.

    \item \textbf{True.}
        Each \julia{eps()/3} is represented almost exactly, and twice the
        value is about $2\varepsilon/3 \approx 1.48 \times 10^{-16}$,
        which is larger than the rounding threshold
        $\varepsilon/2 \approx 1.11 \times 10^{-16}$.
        Adding this quantity to 1 therefore rounds upward to
        $1 + \varepsilon$, so the left-hand side is strictly larger than 1.
        (Confirmed in Julia.)

    \item \textbf{False.}
        Doubling can fall outside the representable range (overflow to
        $\infty$, underflow of subnormals) or introduce rounding: e.g.\
        $2 \cdot x$ for the largest finite number is not finite.

    \item \textbf{False.}
        Both formats represent only finitely many real numbers exactly.

    \item \textbf{True.}
        The rounded result of the multiplication is exact when the product
        fits in the mantissa (e.g.\ powers of two) and is rounded otherwise.
\end{enumerate}

\noindent\textbf{Bonus.}
\begin{enumerate}[itemsep=8pt]
    \item \textbf{False.}
        The true value of $\sqrt{1 + 2\varepsilon}$ is approximately
        $1 + \varepsilon - \varepsilon^2/2$, which rounds to the nearest
        representable number \julia{1 + eps()} (i.e.\ $1 + \varepsilon$),
        not to 1.

    \item \textbf{True.}
        $2 + \varepsilon$ lies exactly at the midpoint between the
        consecutive representable numbers 2 and $2 + 2\varepsilon$
        (the spacing near 2 is $2\varepsilon$), so with round-to-nearest-even
        it rounds to 2. Consequently
        $\log_2(2 + \varepsilon)$ computes $\log_2(2) = 1$ exactly.
        (Confirmed in Julia.)
\end{enumerate}
\fi

\end{document}
